\section{Appendix: Vector Field Environments}
\label{sec:appendix_vector_fields}

Six vector field environments centered at the origin were used to evaluate the multi-robot navigation algorithms. Let $r = \sqrt{x^2 + y^2}$ and $\theta = \arctan2(y, x)$.

\subsection{Field Equations}

\textbf{Vortex:} Pure rotational flow.
\begin{equation}
\mathbf{v}_{\text{vortex}}(x, y) = \begin{bmatrix} -r \sin(\theta) \\ r \cos(\theta) \end{bmatrix}
\end{equation}

\textbf{Sinking Vortex:} Spiraling inward flow.
\begin{equation}
\mathbf{v}_{\text{sink-vortex}}(x, y) = 0.4 \begin{bmatrix} -r \sin(\theta) \\ r \cos(\theta) \end{bmatrix} + 0.15 \begin{bmatrix} -x \\ -y \end{bmatrix}
\end{equation}

\textbf{Sink:} Radial inward flow.
\begin{equation}
\mathbf{v}_{\text{sink}}(x, y) = \begin{bmatrix} -x \\ -y \end{bmatrix}
\end{equation}

\textbf{Source:} Radial outward flow.
\begin{equation}
\mathbf{v}_{\text{source}}(x, y) = \begin{bmatrix} x \\ y \end{bmatrix}
\end{equation}

\textbf{Spewing Vortex:} Spiraling outward flow.
\begin{equation}
\mathbf{v}_{\text{spew-vortex}}(x, y) = -0.4 \begin{bmatrix} -r \sin(\theta) \\ r \cos(\theta) \end{bmatrix} - 0.15 \begin{bmatrix} -x \\ -y \end{bmatrix}
\end{equation}

\textbf{Saddle:} Hyperbolic saddle point rotated $45^\circ$.
\begin{equation}
\mathbf{v}_{\text{saddle}}(x, y) = \mathbf{R}(-\pi/4) \begin{bmatrix} x' \\ -y' \end{bmatrix}, \quad \text{where } \begin{bmatrix} x' \\ y' \end{bmatrix} = \mathbf{R}(\pi/4) \begin{bmatrix} x \\ y \end{bmatrix}
\end{equation}

\subsection{Classification}

\begin{table}[h]
\centering
\begin{tabular}{lccc}
\hline
\textbf{Field Type} & \textbf{Divergence} & \textbf{Curl} & \textbf{Critical Point} \\
\hline
Vortex & 0 & $\neq 0$ & Center \\
Sinking Vortex & $< 0$ & $\neq 0$ & Stable Spiral \\
Sink & $< 0$ & 0 & Stable Node \\
Source & $> 0$ & 0 & Unstable Node \\
Spewing Vortex & $> 0$ & $\neq 0$ & Unstable Spiral \\
Saddle & 0 & 0 & Saddle Point \\
\hline
\end{tabular}
\caption{Topological classification of vector field environments.}
\label{tab:field_summary}
\end{table}
